% !TeX program = xelatex
% 修改标题、汇报人和正文即可开始使用。XeLaTeX 编译两遍。
\documentclass[aspectratio=169]{beamer}
% Keep the original WenKai/Sarasa typography without installing fonts.
% A small Overleaf package can omit fonts/ and use TeX Live's Fandol instead.
\usepackage{fontspec}
\usepackage{xeCJK}
\IfFileExists{fonts/LXGWWenKai-Regular.ttf}{%
  \setCJKmainfont{LXGWWenKai-Regular.ttf}[
    Path=fonts/, BoldFont=LXGWWenKai-Medium.ttf,
    ItalicFont=LXGWWenKai-Regular.ttf, ItalicFeatures={FakeSlant=0.15}]
  \setCJKsansfont{LXGWWenKai-Regular.ttf}[
    Path=fonts/, BoldFont=LXGWWenKai-Medium.ttf,
    ItalicFont=LXGWWenKai-Regular.ttf, ItalicFeatures={FakeSlant=0.15}]
  \setCJKmonofont{SarasaMonoSC-Regular.ttf}[
    Path=fonts/, ItalicFont=SarasaMonoSC-Regular.ttf,
    ItalicFeatures={FakeSlant=0.15}]
}{%
  \setCJKmainfont{FandolSong-Regular.otf}[
    BoldFont=FandolSong-Bold.otf, ItalicFont=FandolKai-Regular.otf]
  \setCJKsansfont{FandolHei-Regular.otf}[
    BoldFont=FandolHei-Bold.otf, ItalicFont=FandolHei-Regular.otf,
    ItalicFeatures={FakeSlant=0.15}]
  \setCJKmonofont{FandolFang-Regular.otf}[
    ItalicFont=FandolFang-Regular.otf, ItalicFeatures={FakeSlant=0.15}]
}
\IfFontExistsTF{Times New Roman}{\setmainfont{Times New Roman}}{%
  \setmainfont{texgyretermes-regular.otf}[
    BoldFont=texgyretermes-bold.otf, ItalicFont=texgyretermes-italic.otf,
    BoldItalicFont=texgyretermes-bolditalic.otf]}
\IfFontExistsTF{Arial}{\setsansfont{Arial}}{%
  \setsansfont{texgyreheros-regular.otf}[
    BoldFont=texgyreheros-bold.otf, ItalicFont=texgyreheros-italic.otf,
    BoldItalicFont=texgyreheros-bolditalic.otf]}
\setmonofont{lmmono10-regular.otf}[
  BoldFont=lmmonolt10-bold.otf, ItalicFont=lmmono10-italic.otf,
  BoldItalicFont=lmmonolt10-boldoblique.otf]

\usetheme{HRBUST}

\title[算法训练与竞赛交流]{哈尔滨理工大学 ACM-ICPC 集训队}
\subtitle{算法训练与竞赛交流}
\author[ACM-ICPC]{汇报人}
\institute[哈理工集训队]{哈尔滨理工大学 ACM-ICPC 集训队}
\date{\today}
\hypersetup{colorlinks=true,linkcolor=.,urlcolor=hrbustnavy}

\begin{document}

\begin{frame}
  \titlepage
  \begin{figure}[htpb]
    \centering
    \hrbustlogo
  \end{figure}
\end{frame}

\begin{frame}
  \tableofcontents
\end{frame}

\section{问题与思路}

\begin{frame}{单源最短路}
  给定一个含 \(n\) 个点、\(m\) 条边的有向图，边权均为非负数。
  求从起点 \(s\) 到每个点的最短距离。
  \begin{block}{距离定义}
    \(\mathrm{dist}[v]\) 表示从 \(s\) 到 \(v\) 的当前最短距离。
    初始令 \(\mathrm{dist}[s]=0\)，其余距离为 \(\infty\)。
  \end{block}
  \begin{itemize}
    \item 维护一个按距离排序的优先队列
    \item 每次取出距离最小的点，尝试更新相邻点的距离
  \end{itemize}
\end{frame}

\begin{frame}{状态更新与适用条件}
  \begin{block}{松弛操作}
    对边 \((u,v)\) 及其权重 \(w\)，更新
    \[
      \mathrm{dist}[v] = \min\bigl(\mathrm{dist}[v],\mathrm{dist}[u]+w\bigr).
    \]
  \end{block}
  \begin{alertblock}{适用条件}
    Dijkstra 算法要求边权非负。距离与边权应使用足够大的整数类型。
  \end{alertblock}
  \begin{exampleblock}{队列处理}
    若弹出的距离与 \(\mathrm{dist}[u]\) 不同，跳过该过期状态。
  \end{exampleblock}
\end{frame}

\section{实现与复杂度}

\begin{frame}[fragile]{核心实现}
  \begin{lstlisting}[language=C++,basicstyle=\ttfamily\footnotesize]
using State = pair<long long, int>;
priority_queue<State, vector<State>, greater<State>> pq;
dist[s] = 0;
pq.push({0, s});
while (!pq.empty()) {
    auto [d, u] = pq.top(); pq.pop();
    if (d != dist[u]) continue;
    for (auto [v, w] : adj[u]) {
        if (d + w < dist[v]) {
            dist[v] = d + w;
            pq.push({dist[v], v});
        }
    }
}
  \end{lstlisting}
\end{frame}

\begin{frame}{复杂度与边界检查}
  \begin{table}
    \centering
    \begin{tabular}{ll}
      \toprule
      项目 & 说明 \\
      \midrule
      时间复杂度 & \(O((n+m)\log(n+m))\)，二叉堆并允许重复入队 \\
      空间复杂度 & \(O(n+m)\)，邻接表与优先队列 \\
      不可达顶点 & 距离保持为 \(\infty\) \\
      整数范围 & 使用 \texttt{long long}，避免距离相加溢出 \\
      \bottomrule
    \end{tabular}
  \end{table}
  \begin{itemize}
    \item 检查重边、零权边以及只有一个点的情况
    \item 按题目要求输出不可达状态
  \end{itemize}
\end{frame}

\section{训练复盘}

\begin{frame}{训练记录}
  \begin{columns}[t]
    \begin{column}{0.48\textwidth}
      \begin{block}{本次训练}
        \begin{itemize}
          \item 题目与关键观察
          \item 正确性证明
          \item 实现中的边界条件
        \end{itemize}
      \end{block}
    \end{column}
    \begin{column}{0.48\textwidth}
      \begin{exampleblock}{后续练习}
        \begin{itemize}
          \item 对照小规模暴力算法检查
          \item 记录错误原因与修正方法
          \item 补充同类题目
        \end{itemize}
      \end{exampleblock}
    \end{column}
  \end{columns}
\end{frame}

\begin{frame}
  \begin{center}
    {\Huge\calligra Thanks!}
  \end{center}
\end{frame}

\end{document}
